\documentclass[11pt]{article}

\usepackage[margin=1in]{geometry}
\usepackage{amsmath, amssymb, bm}
\usepackage{booktabs}
\usepackage{array}
\usepackage{hyperref}
\usepackage{xcolor}

\title{Complete PhiBE-TD3-BC for Continuous-Time Glucose Control}
\author{Hanzhang Yuan}
\date{\today}

\begin{document}
\maketitle

\begin{abstract}
This report summarizes the current progress of a continuous-time offline reinforcement
learning project for type 1 diabetes glucose control. The project starts from the TD3-BC
offline RL baseline in Emerson et al. and adds a physics-informed Bellman equation
(PhiBE) regularization term motivated by the continuous-time Hamilton--Jacobi--Bellman
(HJB) equation. The main methodological update is a complete second-order PhiBE loss
that includes the full Hessian/diffusion contribution, rather than only a first-order drift
term or a diagonal approximation. Experiments are performed on the adult\#1 virtual
patient with training data collected at the original 3-minute interval and evaluation at
1, 3, and 5 minute sampling intervals. The current best stable Complete PhiBE setting is
\(\lambda_{\mathrm{PhiBE}}=0.01\), \(\alpha=1.25\), and 70k training updates. This setting
improves time in range across all tested sampling intervals, but slightly increases time
below range, suggesting that the next research direction should focus on preserving the
TIR improvement while controlling hypoglycemia risk.
\end{abstract}

\section{Background and Objective}

The baseline project follows the offline RL glucose-control framework of Emerson et al.,
where a policy is trained from a fixed replay buffer rather than by interacting with a
patient online. This is important for glucose control because unsafe online exploration
can lead to clinically undesirable insulin dosing.

The current project investigates whether the TD3-BC baseline can be improved by using
continuous-time structure. Instead of treating glucose control only as a discrete-time
Bellman learning problem, we regularize the critic using a residual inspired by the
continuous-time HJB equation. The research question is:

\begin{quote}
Can a continuous-time PhiBE regularizer improve glucose-control performance under
different evaluation sampling intervals while maintaining acceptable hypoglycemia risk?
\end{quote}

\section{Baseline: TD3-BC}

TD3-BC trains an actor \(\pi_\phi(s)\) and two critics \(Q_{\theta_1}(s,a)\),
\(Q_{\theta_2}(s,a)\). For a batch transition
\((s_i,a_i,r_i,s'_i)\), the TD target is
\begin{equation}
y_i
= r_i
+ \gamma
\min_{j=1,2}
Q_{\bar{\theta}_j}
\left(
s'_i,
\pi_{\bar{\phi}}(s'_i) + \epsilon
\right),
\end{equation}
where \(\epsilon\) is clipped target-policy smoothing noise. The critic loss is
\begin{equation}
\mathcal{L}_{\mathrm{TD}}
=
\frac{1}{B}
\sum_{i=1}^{B}
\left[
\left(Q_{\theta_1}(s_i,a_i)-y_i\right)^2
+
\left(Q_{\theta_2}(s_i,a_i)-y_i\right)^2
\right].
\end{equation}

The TD3-BC actor objective balances value maximization and behavioral cloning:
\begin{equation}
\mathcal{L}_{\pi}
=
-
\frac{\alpha}
{\frac{1}{B}\sum_{i=1}^{B}\left|Q_{\theta_1}(s_i,\pi_\phi(s_i))\right|}
\cdot
\frac{1}{B}\sum_{i=1}^{B} Q_{\theta_1}(s_i,\pi_\phi(s_i))
+
\frac{1}{B}\sum_{i=1}^{B}
\left\|\pi_\phi(s_i)-a_i\right\|_2^2.
\end{equation}
Here, \(\alpha\) controls how strongly the actor follows the critic value term relative
to the behavior-cloning term.

\section{Continuous-Time PhiBE Formulation}

We model the underlying glucose dynamics as a stochastic differential equation:
\begin{equation}
d s_t
= b(s_t,a_t)\,dt
+ \sigma(s_t,a_t)\,dW_t,
\end{equation}
where \(b(s_t,a_t)\) is the drift term, \(\sigma(s_t,a_t)\) is the diffusion term, and
\(W_t\) is Brownian motion.

The corresponding continuous-time optimal-control equation is the HJB equation:
\begin{equation}
\beta V^\star(s)
=
\max_a
\left[
r(s,a)
+
b(s,a)^\top \nabla_s V^\star(s)
+
\frac{1}{2}
\sigma(s,a)\sigma(s,a)^\top
:
\nabla_s^2 V^\star(s)
\right].
\end{equation}

In implementation, the critic \(Q_\theta(s,a)\) is used as the value-like function.
For each replay transition, define
\begin{equation}
\Delta s_i = s'_i - s_i,
\qquad
\Delta t = 3 \ \mathrm{min}.
\end{equation}
The empirical drift is approximated as
\begin{equation}
\hat{b}_i = \frac{\Delta s_i}{\Delta t}.
\end{equation}

The Complete PhiBE residual is
\begin{equation}
\mathcal{R}_{\theta}(s_i,a_i)
=
\beta Q_{\theta}(s_i,a_i)
- r_i
- \nabla_s Q_{\theta}(s_i,a_i)^\top \hat{b}_i
- \frac{1}{2\Delta t}
\Delta s_i^\top
\nabla_s^2 Q_{\theta}(s_i,a_i)
\Delta s_i .
\end{equation}

The last term is the full Hessian/diffusion term. It includes cross-state curvature
effects through the quadratic form
\begin{equation}
\Delta s_i^\top
\nabla_s^2 Q_{\theta}(s_i,a_i)
\Delta s_i.
\end{equation}
The implementation computes this quantity using a Hessian-vector product, so it does not
explicitly materialize a full Hessian matrix for every batch sample.

The PhiBE regularization loss is
\begin{equation}
\mathcal{L}_{\mathrm{PhiBE}}
=
\frac{1}{B}
\sum_{i=1}^{B}
\left[
\mathcal{R}_{\theta_1}(s_i,a_i)^2
+
\mathcal{R}_{\theta_2}(s_i,a_i)^2
\right].
\end{equation}

The final critic objective is
\begin{equation}
\mathcal{L}_{\mathrm{critic}}
=
\mathcal{L}_{\mathrm{TD}}
+
\lambda_{\mathrm{PhiBE}}
\mathcal{L}_{\mathrm{PhiBE}}.
\end{equation}

\section{Experimental Setup}

The current experiments use adult\#1 from the simulator. The replay buffer is collected
with the original 3-minute setting used by the baseline protocol. To avoid confounding the
comparison, the training replay construction is kept fixed at 3 minutes. The trained policy
is then evaluated under multiple continuous-time testing intervals:
\[
\Delta t \in \{1,3,5\} \ \mathrm{min}.
\]

The main metrics are:
\begin{itemize}
\item Time in range (TIR): percentage of glucose values in 70--180 mg/dl.
\item Time below range (TBR): percentage of glucose values below 70 mg/dl.
\item Time above range (TAR): percentage of glucose values above 180 mg/dl.
\item Mean blood glucose (Mean BG).
\item Mean insulin action.
\end{itemize}

\section{Current Results}

The current best stable Complete PhiBE configuration is
\[
\lambda_{\mathrm{PhiBE}}=0.01,
\qquad
\alpha=1.25,
\qquad
\mathrm{training\ updates}=70{,}000.
\]

\begin{table}[h]
\centering
\caption{Adult\#1 multi-\(\Delta t\) evaluation.}
\label{tab:main_results}
\begin{tabular}{llrrrrr}
\toprule
Method & \(\Delta t\) & \(N\) & TIR & TBR & TAR & Mean BG \\
\midrule
TD3-BC & 1 min & 9 & 88.71\% & 2.56\% & 8.73\% & 131.7 \\
PID & 1 min & 3 & 85.64\% & 2.00\% & 12.36\% & 137.0 \\
Complete PhiBE & 1 min & 9 & 89.06\% & 4.51\% & 6.43\% & 125.3 \\
\midrule
TD3-BC & 3 min & 9 & 69.26\% & 0.07\% & 30.66\% & 156.8 \\
PID & 3 min & 3 & 61.04\% & 0.04\% & 38.92\% & 164.5 \\
Complete PhiBE & 3 min & 9 & 71.98\% & 0.21\% & 27.81\% & 153.1 \\
\midrule
TD3-BC & 5 min & 9 & 62.29\% & 0.05\% & 37.66\% & 164.9 \\
PID & 5 min & 3 & 53.19\% & 0.00\% & 46.81\% & 174.1 \\
Complete PhiBE & 5 min & 9 & 65.22\% & 0.09\% & 34.68\% & 161.6 \\
\bottomrule
\end{tabular}
\end{table}

\section{Alpha Sensitivity and Stability}

The first Complete PhiBE setting tested was \(\alpha=1.75\). It improved TIR at 3 and
5 minutes, but one training seed became too aggressive and produced high TBR at 1 minute.
To understand this tradeoff, smaller actor weights were evaluated.

\begin{table}[h]
\centering
\caption{Seed-1 alpha sensitivity for Complete PhiBE.}
\label{tab:alpha_sweep}
\begin{tabular}{lrrrrr}
\toprule
Setting & \(\Delta t\) & TIR & TBR & TAR & Mean BG \\
\midrule
\(\alpha=1.25\) & 1 min & 88.74\% & 6.59\% & 4.66\% & 120.0 \\
\(\alpha=1.25\) & 3 min & 75.44\% & 0.38\% & 24.18\% & 148.7 \\
\(\alpha=1.25\) & 5 min & 70.46\% & 0.20\% & 29.34\% & 154.8 \\
\midrule
\(\alpha=1.30\) & 1 min & 89.52\% & 5.31\% & 5.16\% & 122.1 \\
\(\alpha=1.30\) & 3 min & 73.96\% & 0.40\% & 25.65\% & 150.3 \\
\(\alpha=1.30\) & 5 min & 69.86\% & 0.12\% & 30.02\% & 155.9 \\
\midrule
\(\alpha=1.50\) & 1 min & 77.34\% & 20.30\% & 2.36\% & 102.3 \\
\(\alpha=1.50\) & 3 min & 79.21\% & 2.40\% & 18.40\% & 140.1 \\
\(\alpha=1.50\) & 5 min & 75.10\% & 1.52\% & 23.38\% & 146.6 \\
\bottomrule
\end{tabular}
\end{table}

These results show that increasing \(\alpha\) can improve TIR at 3 and 5 minutes, but the
policy becomes more aggressive and may increase hypoglycemia risk at 1 minute. The current
best stable setting is therefore \(\alpha=1.25\), although the TIR improvement is more
conservative than the more aggressive settings.

\section{Current Interpretation and Next Direction}

Complete PhiBE improves TIR across all three evaluation intervals, especially at 3 and
5 minutes, but the improvement comes with slightly higher TBR. The current interpretation
is that the Hessian/diffusion term changes the critic shape in a way that encourages a more
aggressive insulin policy. This is useful for reducing hyperglycemia, but it can also push
glucose too low when actions are applied more frequently.

The next research question is whether the project should focus on reducing the mild TBR
increase while keeping the TIR gain, or whether another direction should be prioritized
first. A natural next experiment is a safety-aware Complete PhiBE objective that increases
the penalty near or below the hypoglycemia threshold, while keeping the second-order
continuous-time regularization.

\section{Reproducibility Notes}

The current codebase keeps the original 3-minute training replay construction. The main
commands are:
\begin{verbatim}
python run_emerson2023_td3bc_reproduction.py \
  --patients adult#1 --train-seeds 0,1,2 --test-seeds 0,1,2 \
  --replay-length 100000 --training-timesteps 100000

python run_stage1_phibe_adult1.py \
  --train-seed 0 --training-timesteps 100000 --num-train-steps 70000 \
  --phibe-mode full_second_order --lambda-phibe 0.01 --alpha 1.25 \
  --save-tag phibe_stage4_full_second_order_lam1e-2_norm_alpha1p25

python eval_multidt_stage2.py \
  --train-seeds 0,1,2 --test-seeds 0,1,2 --dts 1,3,5 \
  --phibe-mode full_second_order --lambda-phibe 0.01 --alpha 1.25
\end{verbatim}

\begin{thebibliography}{9}

\bibitem{emerson2023}
H. Emerson, M. J. Guy, and R. McConville.
\newblock Offline Reinforcement Learning for Safer Blood Glucose Control in People with Type 1 Diabetes.
\newblock \emph{Journal of Biomedical Informatics}, 142:104376, 2023.
\newblock doi: \href{https://doi.org/10.1016/j.jbi.2023.104376}{10.1016/j.jbi.2023.104376}.

\end{thebibliography}

\end{document}
